\chapter{Code Generation on the Legacy Core}\label{r04:ch:codegen}

A legacy application never contains its lattice Boltzmann kernels as source. A Python
generator script describes the method, and the build runs that script to produce C++
files. This chapter follows that path (\cref{r04:fig:codegen-flow}) and marks the one
place where the vector instruction set is decided.

\begin{figure}[htbp]
  \centering
  \resizebox{\linewidth}{!}{% Legacy code generation: from CMake to compiled sweeps (chapter 3).
\begin{tikzpicture}[node distance=5mm and 6mm]
  % configure time
  \node[r04box,text width=40mm] (cmake) {CMake configure\\[1pt]
    \scriptsize\code{waLBerla_generate_target_from_python}\\(NAME, FILE, CODEGEN\_CFG, OUT\_FILES)};
  \node[r04code,right=16mm of cmake,text width=52mm,align=left] (cmd) {python <script>.py\\
    \ -f <output files>\\
    \ -c \{"CODEGEN\_CFG": "aa\_rvv", "WALBERLA\_...": ...\}\\
    PYTHONPATH fixed at configure time};
  \draw[r04arrow] (cmake) -- node[r04label,above]{custom\\command} (cmd);

  % build time: generator script
  \node[r04box,below=9mm of cmake,text width=40mm] (script) {generator script\\[1pt]
    \scriptsize\code{with CodeGeneration() as ctx:}\\\code{ctx.config} = CODEGEN\_CFG};
  \draw[r04arrow] (cmd.south) |- ($(script.north)+(0,3mm)$) -- (script.north);
  \node[r04label,anchor=west] at ($(cmd.south)+(1mm,-3mm)$) {build time: runs when the script changes};

  \node[r04box,below=of script,text width=40mm] (bridge) {\code{lbmpy_walberla},
    \code{pystencils_walberla}\\[1pt]\scriptsize\code{generate_lbm_package},
    \code{generate_sweep},\\\code{config_from_context}};
  \draw[r04arrow] (script) -- (bridge);

  \node[r04box,right=16mm of bridge,text width=52mm] (ps) {\ps{}~1.4 and lbmpy~1.4\\[1pt]
    \scriptsize collision rule $\to$ AST $\to$ \code{vectorize()} $\to$ C code\\
    with the instruction set's templates (chapter 4)};
  \draw[r04arrow] (bridge) -- node[r04label,above]{kernels} (ps);

  \node[r04box,below=of bridge,text width=40mm] (jinja) {Jinja2 templates\\[1pt]
    \scriptsize\file{templates/*.tmpl.h}, \file{*.tmpl.cpp}\\filters: \code{generate_call}, \dots};
  \draw[r04arrow] (bridge) -- (jinja);
  \draw[r04arrow] (ps.south) |- node[r04label,pos=0.75,below]{kernel C code} (jinja.east);

  \node[r04box,below=of jinja,text width=40mm] (files) {generated C++ files\\[1pt]
    \scriptsize StorageSpecification, SweepCollection,\\BoundaryCollection, \dots};
  \draw[r04arrow] (jinja) -- (files);

  \node[r04box,right=16mm of files,text width=52mm] (lib) {static library \code{<NAME>}\\[1pt]
    \scriptsize compiled with the build's flags (\code{-march=...})\\linked into the application};
  \draw[r04arrow] (files) -- (lib);
\end{tikzpicture}
}
  \caption[Legacy code generation, from CMake to compiled sweeps]{Legacy code generation.
    At configure time CMake writes a custom command per generator target; at build time
    the command runs the generator script, which calls the waLBerla bridge packages,
    which call \ps{}~1.4 for the kernels and Jinja2 for the C++ files around them.}
  \label{r04:fig:codegen-flow}
\end{figure}

\section{The CMake side}\label{r04:sec:cmake}

An application declares each generated library with
\code{waLBerla_generate_target_from_python} (\file{cmake/waLBerlaHelperFunctions.cmake}).
waLBerla-wind's benchmark app (\cref{r04:sec:wind-app}) does it once per variant:
\begin{r04cmakecode}
foreach( variant ${WIND_BENCHMARK_VARIANTS} )          # e.g. pull_none;aa_rvv
    waLBerla_generate_target_from_python(
            NAME TurbineBenchmark_cpu_uniform_${variant}_Generation
            FILE LBMGeneration.py
            CODEGEN_CFG ${variant}
            OUT_FILES waLBerlaWind_SweepCollection.cpp waLBerlaWind_SweepCollection.h ... )
    add_executable( TurbineBenchmark_cpu_uniform_${variant} main.cpp )
    ...
endforeach()
\end{r04cmakecode}
The function does three things:
\begin{enumerate}
  \item It collects CMake variables into a JSON string: \code{CODEGEN_CFG} and the
    switches the script may need, such as \code{WALBERLA_DOUBLE_ACCURACY},
    \code{WALBERLA_BUILD_WITH_OPENMP} and \code{WALBERLA_OPTIMIZE_FOR_LOCALHOST}.
  \item It adds a custom command whose outputs are the \code{OUT_FILES} (in a folder
    named after \code{CODEGEN_CFG}) and which runs
    \code{python <FILE> -f <outputs> -c <json>} with
    \code{PYTHONPATH=<walberla>/python:$ENV{PYTHONPATH}}.
  \item It adds a static library of the generated sources, which the application links.
\end{enumerate}
Two details of this function matter in practice; both are pitfalls in appendix~B.
\code{$ENV{PYTHONPATH}} is expanded when CMake \emph{configures}, so the \ps{} that the
build uses is the one on \code{PYTHONPATH} at configure time; this is how the
\ps{}~1.4 checkout of the submodule is used without installing it. And the custom
command \code{DEPENDS} only on the generator script: a change to \ps{} or to the
bridge packages does not regenerate anything.

\section{The generator script}

A generator script opens a \code{CodeGeneration} context
(\file{python/pystencils_walberla/cmake_integration.py}), which parses the two
arguments and writes the files when the \code{with} block ends:
\begin{pycode}
with CodeGeneration() as ctx:
    streaming_pattern, _, instruction_set = ctx.config.partition('_')   # "aa_rvv" -> "aa", "rvv"
    data_type = 'double' if ctx.double_accuracy else 'float32'
    ...
    cpu_vectorise_info = {'nontemporal': True}
    if instruction_set != 'default':
        cpu_vectorise_info['instruction_set'] = None if instruction_set == 'none' else instruction_set
    generate_lbm_package(ctx, name="waLBerlaWind_", collision_rule=collision_rule, ...,
                         cpu_vectorize_info=cpu_vectorise_info)
\end{pycode}
\code{ctx.config} is the \code{CODEGEN_CFG} string; it is the only channel through
which one script can produce several variants. waLBerla's \code{UniformGridCPU}
encodes stencil, streaming pattern and collision model in it
(\code{d3q19_aa_srt}); our variants append the instruction set
(\cref{r04:ch:apps}).

\section{Where the instruction set is decided}\label{r04:sec:where-isa}

Every kernel of the bridge packages is created with a \ps{} configuration from
\code{config_from_context} (\file{python/pystencils_walberla/utility.py}). It merges
what the script passed in \code{cpu_vectorize_info} with defaults:
\begin{center}\small
\begin{tabularx}{\linewidth}{@{}l>{\raggedright\arraybackslash}X@{}}
  \toprule
  \textbf{key} & \textbf{default}\\
  \midrule
  \code{instruction_set} & \code{get_vectorize_instruction_set(ctx)}: \code{None} (no
    explicit vectors) unless waLBerla is configured with
    \code{WALBERLA_OPTIMIZE_FOR_LOCALHOST=ON}; then the last entry of
    \code{get_supported_instruction_sets()} of the build machine, or \code{'sse'} if
    detection fails\\
  \code{assume_inner_stride_one} & \code{True}\\
  \code{assume_aligned} & \code{False}\\
  \code{nontemporal} & \code{False}\\
  \code{assume_sufficient_line_padding} & \code{False}\\
  \bottomrule
\end{tabularx}
\end{center}
A key the script sets wins over the default. So there are two ways to get vector code:
the script names an instruction set, or the whole build opts in with
\code{OPTIMIZE_FOR_LOCALHOST}, which then applies to \emph{every} kernel that does not
name one. On the K1, detection finds \isa{rvv} (\cref{r04:sec:detect}). This project
names the instruction set per variant and leaves \code{OPTIMIZE_FOR_LOCALHOST} off;
\cref{r04:sec:localhost} explains why it has to be off for waLBerla-wind.

\section{What the LBM package generates}

\code{generate_lbm_package} (\file{python/lbmpy_walberla/walberla_lbm_package.py}) writes
three groups of C++ classes for the module \file{src/lbm_generated}:
\begin{description}[style=nextline]
  \item[StorageSpecification] the PDF field's stencil and layout, and the generated
    pack and unpack kernels for the halo exchange. For in-place patterns only the
    directions that the next step reads are sent.
  \item[SweepCollection] the time-step kernels: \code{streamCollide}, \code{collide},
    \code{stream}, \code{streamOnlyNoAdvancement}, plus \code{initialise} (PDFs from
    density and velocity) and \code{calculateMacroscopicParameters}. Only the first four
    are vectorized; the setter and getter are generated without vectors
    (\code{config_unoptimized} in \file{sweep_collection.py}).
  \item[BoundaryCollection] the boundary conditions, one generated class each.
\end{description}
For in-place streaming every kernel exists twice, as an even and an odd step. The
generated per-block sweep reads the parity from the PDF field and advances it:
\begin{ccode}
void streamCollide(IBlock * block) {                       // waLBerlaWind_SweepCollection.h
   ...
   uint8_t timestep = pdfs->advanceTimestep();
   streamCollide(density, force, pdfs, velocity, omega, timestep, ghost_layers);
}                                                          // picks the Even or Odd kernel
\end{ccode}
For in-place patterns the package also writes \term{block-wise} kernels,
\code{blockStreamCollide} and \code{ghostLayerPropagation}, which process all blocks of
a refinement level in one call through arrays of field pointers
(\code{_data_pdfs_dp}). They are compiled into every in-place variant even when the
application uses the per-block sweeps, as waLBerla-wind's benchmark does; this is why a
bug in them stopped the K1 build (\cref{r04:sec:blockwise-bug}).

\section{Around the kernels: the templates}

The C++ code around each kernel comes from Jinja2 templates
(\file{python/lbmpy_walberla/templates/}, \file{python/pystencils_walberla/templates/})
and their filters (\file{python/pystencils_walberla/jinja_filters.py}). The filter
\code{generate_call} declares a local variable for every parameter of the kernel
function (data pointers, sizes, strides) from the waLBerla field objects, then calls
the kernel. For a kernel over a whole field, for example:
\begin{ccode}
float * RESTRICT  _data_pdfs = pdfs->dataAt(-ghost_layers, -ghost_layers, -ghost_layers, 0);
const int64_t _size_pdfs_0   = int64_t(int64_c(pdfs->xSize()) + 2*ghost_layers);
const int64_t _stride_pdfs_1 = int64_t(pdfs->yStride());
...
internal_..._streamCollideEven::..._streamCollideEven(_data_density, _data_force, _data_pdfs, ...);
\end{ccode}
The kernel's parameter list is decided by \ps{}: it contains exactly the sizes and
strides the kernel's code uses. The template code has to provide all of them. The
block-wise call is the one place where the template does not ask the kernel but
assumes which sizes it needs, and that assumption broke with vector code
(\cref{r04:sec:blockwise-bug}).
